% Conversão fiel da participação como cliente Hoortech; origem imutável RELATORIO_ATUAL.odt.
\hypertarget{__RefHeading___Toc4802_37900007991}{}
\hypertarget{_Toc456271431}{}

% SOURCE text-h-0024
\refstepcounter{section}
\section*{6.1 Introdução}
\label{cliente-hoortech:text-h-0024}
\addcontentsline{toc}{section}{6.1 Introdução}

% SOURCE text-p-0658
\originaltab{}No primeiro semestre de 2023, participei da Hackatona de Engenharia de Software, cujo tema era o desenvolvimento de uma Inteligência Artificial voltada para o bem social. Eu e minha equipe decidimos utilizar um modelo de rede neural convolucional que reconhece a Língua Brasileira de Sinais (Libras) através da \textit{webcam} do usuário para desenvolver uma plataforma para ser utilizada durante reuniões e apresentações. O software processa as imagens capturadas pela câmera, reconhece os sinais em Libras e os converte em texto exibido na plataforma, facilitando a comunicação entre pessoas surdas/mudas e indivíduos que não compreendem a Língua Brasileira de Sinais.\par

% SOURCE text-p-0660
\originaltab{}Após sermos premiados com a segunda colocação na competição, fomos convidados a participar da edição de 2024/2 da AGES, onde nosso projeto foi desenvolvido ao longo de um semestre pelos alunos de Engenharia de Software e Sistemas de Informação. Durante esse período, a equipe aplicou metodologias ágeis, conduzindo \textit{sprints} regulares para entregas incrementais, garantindo uma evolução contínua e estruturada da solução.\par

% FIGURE draw-image-0006; source text-p-0663; image provenance pending human identification.
\memorialfigure{imagens/original/asset-0017.jpg}{\textbf{Figura 33 –} Foto do time - CLIENTE AGES}{cliente-hoortech:draw-image-0006}{Fonte: pendente de identificação}{0.92}

